\section{Discussion}
\label{sec:discussion}

\subsection{What the Refutation Does and Does Not Show}
The results reject three specific propositions on one dataset: that this temporal KAN module improves accuracy over a budget-matched plain KAN, that fusing embedding shift with interval width yields a better reliability signal than width alone, and that the resulting score can identify forecasts that will fail. They do not show that KANs are unsuitable for climate forecasting, that shift detection is useless in general, or that selective prediction does not work. The plain KAN is a competitive model here, and width-based selection works well; what fails is the specific composition we proposed.

The generality of the refutation is nonetheless asymmetric with respect to the positive claim it replaces. Had \method{} won, a single station over a single record would have been weak evidence for a general architectural advantage. Because it loses, and loses to the ablation designed to isolate its central component, the burden shifts: a design that cannot beat its own control on a long, clean, well-understood series has not earned the presumption that it would succeed elsewhere. We report the result at this scope and state the extension to station-generalization data as unfinished work rather than as a defence.

\subsection{Why the Fusion Degrades a Working Signal}
The most instructive negative result is that the fused score is worse than one of its own inputs. Interval width carries real information about difficulty, as the width-only ablation demonstrates. The embedding-shift statistic carries little, as the shift-only ablation demonstrates. Combining them at the frozen equal weighting therefore dilutes the informative signal with an uninformative one, and the AURC penalty grows with horizon precisely where the width signal is strongest.

This suggests a methodological point that outlasts the specific design. The fusion weights were frozen at equal values because freezing was treated as sufficient protection against post-hoc tuning. Freezing an arbitrary value is not the same as choosing a defensible one: it prevents the weights from being fitted to the test set, but it does nothing to ensure they are reasonable, and it can lock in a choice that is worse than either endpoint. The correct discipline is to select such constants on the validation partition, which is legitimate and was available to us, and to freeze the selection procedure rather than the resulting number. We did not do this, and the consequence is a component that is worse than its own ablation. We report the frozen result rather than retuning, because retuning after seeing these tables would convert a pre-registered test into an exploratory one; a validation-selected weighting is the appropriate subject of a subsequent, separately pre-registered study.

\subsection{Marginal Coverage Is Not Evidence of Method Quality}
Conformal marginal coverage was the framework's most presentable number and its least meaningful one. Split conformal delivers it for any predictor under exchangeability, so reporting it as a property of \method{} would attribute to the architecture something contributed entirely by the wrapper. The joint-coverage collapse, from $0.894$ to $0.103$ across the horizon range, is the more decision-relevant quantity and points the other way. We note this because marginal coverage tables are common in the trustworthy-forecasting literature and are easy to read as validation of a method rather than of a calibration procedure.

\subsection{Interpretability versus Reliability}
KAN functions can provide inspectable nonlinear relationships, but interpretability and reliability are distinct properties, and this study separates them more sharply than intended. Any interpretability analysis of \method{} now describes a representation that is less accurate than a plain KAN at every horizon beyond one day, so the inspected functions cannot be presented as explaining superior behaviour. Stable functions may still aid scientific inspection; they are simply not evidence of forecast reliability, and the reliability measurements reported here are the relevant test of that.

\subsection{On Finding a Defect Mid-Study}
The mean-pooled readout described in Section~\ref{sec:correction} averaged encoder states over the full history window and discarded recency, leaving the model behind persistence. It was found because a pre-registered baseline comparison made the failure impossible to overlook, not because of a code review. Two aspects are worth recording. First, the defect was silent: training converged, losses looked reasonable, and every artifact validated. Only the comparison against a trivial baseline exposed it. Second, correcting it changed the code fingerprint and invalidated every completed run, which forced a full re-execution; the audit refuses to aggregate across fingerprints precisely so that this cost cannot be avoided by mixing old and new results. The correction improved the model substantially and it still loses, which is the finding reported here.

\subsection{Limitations}
Several limitations bound these conclusions. The headline evidence is a single long station record, so the results establish failure on CET rather than universal ranking; the station-generalization campaign that would test breadth is specified and not yet executed, and the multivariate protocols are frozen but unreported for lack of data access. Conformal guarantees depend on exchangeability assumptions that temporal dependence and nonstationarity strain, which is one reason joint coverage degrades. Mahalanobis-type embedding distance captures only one form of representation shift, so the shift component's poor showing is evidence against this operationalization rather than against shift-aware reliability as an idea. Selective prediction is only meaningful where abstention or escalation is operationally available. Finally, the comparison family was fixed at two comparators; a wider family would give a fuller picture of where the design sits, at the cost of multiplicity corrections that were not pre-specified.

\subsection{Reproducibility and Scientific Integrity}
Every headline number is generated from immutable, checksum-bound experiment outputs by a script that refuses to build a table from runs produced under superseded code. The claim--evidence matrix records outcomes against the evidence each claim required, including the rows now marked refuted. Failed runs, the defect and its correction, and the datasets and horizons on which the method underperforms all remain visible. This apparatus did not make the project succeed; it made it possible to know that it did not, which is the more useful of the two outcomes to be able to distinguish.
